\section{Tabel dan Algoritma Operator Precedence}
\label{sec:ch11-algoritma}

\subsection{Algoritma}

Parser operator precedence menggunakan stack yang menyimpan terminal
dan nonterminal. Pada setiap langkah:
\begin{enumerate}[label=\arabic*.,leftmargin=*]
  \item cari terminal teratas pada stack, sebut \(a\);
  \item bandingkan \(a\) dengan token saat ini \(b\);
  \item jika \(a <\cdot b\), shift \(b\);
  \item jika \(a \doteq b\), shift atau tangani kurung;
  \item jika \(a \cdot> b\), reduce handle teratas;
  \item jika sel kosong, laporkan kesalahan.
\end{enumerate}

Handle adalah substring pada stack yang cocok dengan ruas kanan
aturan. Setelah reduce, parser membandingkan kembali terminal
teratas dengan token saat ini.

\subsection{Contoh Penelusuran}

Untuk masukan \(\code{id} + \code{id} * \code{id}\):
\begin{enumerate}[label=\arabic*.,leftmargin=*]
  \item shift \code{id}, reduce ke \(F\), lalu \(T\), lalu \(E\);
  \item shift \(+\);
  \item shift \code{id}, reduce ke \(F\), lalu \(T\);
  \item bandingkan \(+\) dengan \(*\): \(+ <\cdot *\), shift \(*\);
  \item shift \code{id}, reduce ke \(F\);
  \item bandingkan \(*\) dengan \(\$\): \(* \cdot> \$\), reduce
        \(T \to T * F\);
  \item bandingkan \(+\) dengan \(\$\): \(+ \cdot> \$\), reduce
        \(E \to E + T\);
  \item terima.
\end{enumerate}

\subsection{Implementasi Sederhana}

Kalkulator shift-reduce pada \file{shift\_reduce.c} mengilustrasikan
stack nilai dan operator. Program ini bukan parser Subset-C lengkap,
tetapi menunjukkan mekanisme shift dan reduce. Potongan fungsinya
ditunjukkan pada \lstref{lst:ch11-reduce}.

\includecode[
  language=C,
  firstline=27,
  lastline=40,
  firstnumber=27,
  caption={Fungsi \code{reduce\_top} pada \file{shift\_reduce.c}},
  label={lst:ch11-reduce}
]{../../code/teaching/shift_reduce.c}

Program membaca ekspresi digit tunggal, mendorong nilai dan operator,
lalu mengurangi sesuai precedence. Keluaran program menunjukkan
urutan shift dan reduce, sehingga mahasiswa dapat membandingkannya
dengan penelusuran manual.

\subsection{Deteksi Kesalahan}

Kesalahan terdeteksi ketika relasi precedence tidak terdefinisi.
Sebagai contoh, pada masukan \(\code{id} + + \code{id}\), setelah
shift \(+\) kedua, tabel tidak memiliki relasi yang valid antara
\(+\) dan \(+\). Parser melaporkan kesalahan dan dapat melakukan
pemulihan dengan melewati token sampai titik aman.

\begin{warningbox}
Operator precedence tidak menangani semua konflik. Jika grammar
memiliki produksi \(\varepsilon\) atau dua nonterminal bersebelahan,
tabel precedence tidak dapat dibangun. Gunakan parser LR untuk
grammar yang lebih umum.
\end{warningbox}
